%% main_rtai_artigo.tex, based on abtex2-modelo-artigo.tex, v-1.9.6 laurocesar
%% Copyright 2012-2016 by abnTeX2 group at http://www.abntex.net.br/ 

% ------------------------------------------------------------------------
% ------------------------------------------------------------------------
% abnTeX2: Modelo de Artigo Acadêmico em conformidade com
% ABNT NBR 6022:2003: Informação e documentação - Artigo em publicação 
% periódica científica impressa - Apresentação
% ------------------------------------------------------------------------
% ------------------------------------------------------------------------

\documentclass[
% -- opções da classe memoir --
article,			% indica que é um artigo acadêmico
11pt,				% tamanho da fonte
twoside,			% para impressão apenas no recto. Oposto a twoside
a4paper,			% tamanho do papel
% -- opções da classe abntex2 --
%chapter=TITLE,		% títulos de capítulos convertidos em letras maiúsculas
%section=TITLE,		% títulos de seções convertidos em letras maiúsculas
%subsection=TITLE,	% títulos de subseções convertidos em letras maiúsculas
%subsubsection=TITLE % títulos de subsubseções convertidos em letras maiúsculas
% -- opções do pacote babel --
twocolumn,          % texto em duas colunas
english,			% idioma adicional para hifenização
brazil,				% o último idioma é o principal do documento
sumario=tradicional
]{abntex2}

\usepackage{sty/packs}

% ----
% Se necessário, inclua outros pacotes aqui

% ---

% ---
% Informações de dados para CAPA e FOLHA DE ROSTO
% ---
\titulo{Modelo de Artigo Original para publicação na \emph{Revista Tecnologia e Arquitetura da Informação}}

% --- Para envio da primeira submissão à Revista, deixe os dados dos autores em branco
\autor{Bruno Carvalho C. Souza\thanks{\href{mailto:bruno.souza@cpai.unb.br}{bruno.souza@cpai.unb.br}} \and Lauro César Araujo\thanks{\href{mailto:laurocesar@laurocesar.com}{laurocesar@laurocesar.com}}}

\local{Brasil}
\data{2016.04.22, v-1.0}
% ---

% ----
% Início do documento
% ----
\begin{document}
    
    % Seleciona o idioma do documento (conforme pacotes do babel)
    %\selectlanguage{english}
    \selectlanguage{brazil}
    
    % Retira espaço extra obsoleto entre as frases.
    \frenchspacing 
    
    % ----------------------------------------------------------
    % ELEMENTOS PRÉ-TEXTUAIS
    % ----------------------------------------------------------
    
    \twocolumn[    		% INICIO DE ARTIGO EM DUAS COLUNAS
    
    % página de titulo
    \maketitle
    
    % resumo em português
    \begin{resumoumacoluna}
        Conforme a ABNT NBR 6022:2003, o resumo é elemento obrigatório, constituído de uma sequência de frases concisas e objetivas e não de uma simples enumeração de tópicos, não ultrapassando 250 palavras, seguido, logo abaixo, das palavras representativas do conteúdo do trabalho, isto é, palavras-chave e/ou descritores, conforme a NBR 6028. (\ldots) As palavras-chave devem figurar logo abaixo do resumo, antecedidas da expressão Palavras-chave:, separadas entre si por ponto e finalizadas também por ponto.
        
        \noindent
        \textbf{Palavras-chave}: informação. lógica. tecnologia. arquitetura da informação.
    \end{resumoumacoluna}
    
    % resumo em português
    \renewcommand{\resumoname}{Abstract}
    \begin{resumoumacoluna}
        \begin{otherlanguage*}{english}
            According to ABNT NBR 6022:2003, an abstract in foreign language is a back
            matter mandatory element.
            
            \noindent
            \textbf{Keywords}: information. Logic. Technology. Architecture of Information.
        \end{otherlanguage*}  
    \end{resumoumacoluna}
    
    \vspace{\onelineskip}%
    
    ]  				% FIM DE ARTIGO EM DUAS COLUNAS
    % ---
    
    
    % ----------------------------------------------------------
    % ELEMENTOS TEXTUAIS
    % ----------------------------------------------------------
    \textual
    
    % ----------------------------------------------------------
    % Introdução
    % ----------------------------------------------------------
    \section*{Introdução}
    \addcontentsline{toc}{section}{Introdução}
    
    Este documento e seu código-fonte são exemplos de referência de uso da classe \textsf{abntex2} e do pacote \textsf{abntex2cite}. O documento exemplifica a elaboração de publicação periódica científica impressa produzida conforme a ABNT NBR 6022:2003 \emph{Informação e documentação - Artigo em publicação periódica científica impressa - Apresentação}.
    
    A expressão ``Modelo canônico'' é utilizada para indicar que \abnTeX\ não é modelo específico de nenhuma universidade ou instituição, mas que implementa tão somente os requisitos das normas da ABNT. Uma lista completa das normas observadas pelo \abnTeX\ é apresentada em \citeonline{abntex2classe}.
    
    Sinta-se convidado a participar do projeto \abnTeX! Acesse o site do projeto em \url{http://www.abntex.net.br/}. Também fique livre para conhecer, estudar, alterar e redistribuir o trabalho do \abnTeX, desde que os arquivos modificados tenham seus nomes alterados e que os créditos sejam dados aos autores originais, nos termos da ``The \LaTeX\ Project Public License''\footnote{\url{http://www.latex-project.org/lppl.txt}}.
    
    Encorajamos que sejam realizadas customizações específicas deste documento. Porém, recomendamos que ao invés de se alterar diretamente os arquivos do \abnTeX, distribua-se arquivos com as respectivas customizações. Isso permite que futuras versões do \abnTeX~não se tornem automaticamente incompatíveis com as customizações promovidas. Consulte \citeonline{abntex2-wiki-como-customizar} par mais informações.
    
    Este exemplo deve ser utilizado como complemento do manual da classe \textsf{abntex2} \cite{abntex2classe}, dos manuais do pacote \textsf{abntex2cite} \cite{abntex2cite,abntex2cite-alf} e do manual da classe \textsf{memoir} \cite{memoir}. Consulte o \citeonline{abntex2modelo} para obter exemplos e informações adicionais de uso de \abnTeX\ e de \LaTeX.
    
    % ----------------------------------------------------------
    % Seção de explicações
    % ----------------------------------------------------------
    \section{Exemplos de comandos}
    
    \subsection{Margens}
    
    A norma ABNT NBR 6022:2003 não estabelece uma margem específica a ser utilizada no artigo científico. 
    
    \subsection{Duas colunas}
    
    É comum que artigos científicos sejam escritos em duas colunas.
    
    É possível indicar pontos do texto que se deseja manter em apenas uma coluna, geralmente o título e os resumos. Os resumos em única coluna em documentos com a opção \texttt{twocolumn} devem ser escritos no ambiente \texttt{resumoumacoluna}.
    
    \subsection{Recuo do ambiente \texttt{citacao}}
    
    Na produção de artigos (opção \texttt{article}), pode ser útil alterar o recuo do ambiente \texttt{citacao}.
    
    Quando um documento é produzido com a opção \texttt{twocolumn}, a classe \textsf{abntex2} automaticamente altera o recuo padrão de 4 cm, definido pela ABNT NBR 10520:2002 seção 5.3, para 1.8 cm.
    
    \section{Mais exemplos no Modelo Canônico de Trabalhos Acadêmicos}
    
    Este modelo de artigo é limitado em número de exemplos de comandos, pois são apresentados exclusivamente comandos diretamente relacionados com a produção de artigos.
    
    Para exemplos adicionais de \abnTeX\ e \LaTeX, como inclusão de figuras, fórmulas matemáticas, citações, e outros, consulte o documento \citeonline{abntex2modelo}.
    
    \section{Consulte o manual da classe \textsf{abntex2}}
    
    Consulte o manual da classe \textsf{abntex2} \cite{abntex2classe} para uma referência completa das macros e ambientes disponíveis.
    
    \lipsum[5]
    
    \lipsum[4]
    
    \lipsum[7]
    
    \lipsum[3]
    
    % ---
    % Finaliza a parte no bookmark do PDF, para que se inicie o bookmark na raiz
    % ---
    \bookmarksetup{startatroot}% 
    % ---
    
    % ---
    % Conclusão
    % ---
    \section*{Considerações finais}
    \addcontentsline{toc}{section}{Considerações finais}
    
    \lipsum[1]
    
    \begin{citacao}
        \lipsum[2]
    \end{citacao}
    
    \lipsum[3]
    
    % ----------------------------------------------------------
    % ELEMENTOS PÓS-TEXTUAIS
    % ----------------------------------------------------------
    \postextual                 % Os elementos pós-textuais estão desabilitados nesse modelo
    
    % ----------------------------------------------------------
    % Referências bibliográficas
    % ----------------------------------------------------------
    \bibliography{rtai}
    
    % ----------------------------------------------------------
    % Glossário
    % ----------------------------------------------------------
    %\glossary
    
    % ----------------------------------------------------------
    % Apêndices
    % ----------------------------------------------------------
    \vspace{1cm}
    \begin{apendicesenv}
        
        % ----------------------------------------------------------
        \chapter{Nullam elementum urna vel imperdiet sodales elit ipsum pharetra ligula
            ac pretium ante justo a nulla curabitur tristique arcu eu metus}
        % ----------------------------------------------------------
        \lipsum[55-57]
        
    \end{apendicesenv}
    % ---
    
    % ----------------------------------------------------------
    % Anexos
    % ----------------------------------------------------------
    \cftinserthook{toc}{AAA}
    \vspace{1cm}
    \begin{anexosenv}
        
        % ---
        \chapter{Cras non urna sed feugiat cum sociis natoque penatibus et magnis dis
            parturient montes nascetur ridiculus mus}
        % ---
        
        \lipsum[31]
        
    \end{anexosenv}
\end{document}
